\chapter{Branches, tags and going back}
\label{ch:history}

\section{Branches}\label{branches}

A branch is a separate line of work, which lets you try something without touching your working version.

\begin{lstlisting}[language=bash]
git switch -c idea       # make a branch and move to it
# ... work and commit ...
git switch main          # go back
git merge idea           # bring the work in
\end{lstlisting}

\textbf{What you should see}

\begin{lstlisting}[language=text]
$ git switch -c idea
Switched to a new branch 'idea'

$ git commit -m "Add the median"
[idea 9b0de38] Add the median
 1 file changed, 1 insertion(+)

$ git switch main
Switched to branch 'main'

$ git merge idea
Updating 3545826..9b0de38
Fast-forward
 results.txt | 1 +
 1 file changed, 1 insertion(+)
\end{lstlisting}

If two branches changed the same lines, Git asks you to choose. Open the file, keep what you want, delete the markers, then \texttt{git\ add} and \texttt{git\ commit}.

Here is a real conflict, where one branch set the mean to 13 and \texttt{main} set it to 14. The file then holds both versions between markers:

\textbf{What you should see}

\begin{lstlisting}[language=text]
$ git merge left
Auto-merging results.txt
CONFLICT (content): Merge conflict in results.txt
Automatic merge failed; fix conflicts and then commit the result.

$ cat results.txt
<<<<<<< HEAD
The mean is 14.
=======
The mean is 13.
>>>>>>> left
The median is 11.
\end{lstlisting}

\section{Tags}\label{tags}

A tag gives a version a permanent name, such as what you sent to a journal.

\begin{lstlisting}[language=bash]
git tag -a submitted-2026-11-01 -m "Sent to the journal"
git push origin submitted-2026-11-01
\end{lstlisting}

\textbf{What you should see}

\begin{lstlisting}[language=text]
$ git tag
submitted-2026-11-01

$ git show --no-patch submitted-2026-11-01
tag submitted-2026-11-01
Tagger: Your Name <you@example.org>
Date:   Sun Nov 1 16:26:36 2026 +0000

Sent to the journal

commit 9b0de3816b57169335da0ad020b08419b9cec876
Author: Your Name <you@example.org>
Date:   Sun Nov 1 16:26:35 2026 +0000

    Add the median
\end{lstlisting}

Once pushed, tags are listed on the project\textquotesingle s page on GitHub, each with its commit and a download of the files at that point. These are the tags of my pamphlet project:

\begin{figure}
\centering
\includegraphics[width=\linewidth]{img/github-tags.png}
\caption{The tags of a public repository on GitHub.}
\end{figure}

\section{Going back}\label{going-back}

{\def\LTcaptype{none} % do not increment counter
\begin{longtable}[]{@{}
  >{\raggedright\arraybackslash}p{(\linewidth - 2\tabcolsep) * \real{0.5000}}
  >{\raggedright\arraybackslash}p{(\linewidth - 2\tabcolsep) * \real{0.5000}}@{}}
\toprule\noalign{}
\begin{minipage}[b]{\linewidth}\raggedright
You want to
\end{minipage} & \begin{minipage}[b]{\linewidth}\raggedright
Use
\end{minipage} \\
\midrule\noalign{}
\endhead
\bottomrule\noalign{}
\endlastfoot
See an old version & \texttt{git\ switch\ -\/-detach\ \textless{}tag\ or\ commit\textgreater{}}, then \texttt{git\ switch\ main} to return \\
Undo a commit, keeping the history & \texttt{git\ revert\ \textless{}commit\textgreater{}} \\
Get one old file back & \texttt{git\ restore\ -\/-source\ \textless{}commit\textgreater{}\ file.txt} \\
Discard edits you haven\textquotesingle t committed & \texttt{git\ restore\ file.txt}, which can\textquotesingle t be undone \\
Find which commit broke something & \texttt{git\ bisect} \\
\end{longtable}
}

\textbf{What you should see: looking at an old version}

\begin{lstlisting}[language=text]
$ git switch --detach submitted-2026-11-01
HEAD is now at 9b0de38 Add the median

$ cat results.txt
The mean is 12.
The median is 11.

$ git switch main
Previous HEAD position was 9b0de38 Add the median
Switched to branch 'main'
\end{lstlisting}

Prefer \texttt{revert} over rewriting history. Commands such as \texttt{git\ reset\ -\/-hard} and \texttt{git\ push\ -\/-force} can destroy work, so check what they do before using them.

\section{A sensible rule for versions}\label{a-sensible-rule-for-versions}

Keep one working copy of each file under Git, and stop putting versions in file names. For data, which is too large for Git, use dated read-only folders with recorded checksums, as described in the \href{../backups/}{backups chapter}.

